\documentclass[11pt]{article}
\usepackage[a4paper,margin=1in]{geometry}
\usepackage{amsmath,amssymb,amsthm,mathtools}
\usepackage{hyperref}

\title{Existence of \(\Omega_p\): An Ultrapower Construction from the Prime Sequence}
\author{Author Name}
\date{\today}

\theoremstyle{definition}
\newtheorem{definition}{Definition}[section]
\newtheorem{remark}[definition]{Remark}
\newtheorem{conjecture}[definition]{Conjecture}

\theoremstyle{plain}
\newtheorem{lemma}[definition]{Lemma}
\newtheorem{theorem}[definition]{Theorem}
\newtheorem{proposition}[definition]{Proposition}

\begin{document}
\maketitle

\begin{abstract}
We define an element \(\Omega_p\) in a nonstandard extension of \(\mathbb N\) as the ultraproduct class of the prime sequence.
We prove existence and basic infinitude properties.
We explicitly separate proved statements from conjectural extensions.
\end{abstract}

\section{Introduction}
This paper establishes the existence of an infinite hypernatural generated by the prime sequence via ultrapower.
No claim is made here about proving the Riemann Hypothesis.
The goal is foundational: existence and first-order properties only.

\section{State of the Art}
Ultrapowers and nonstandard models of arithmetic are classical.
Given a nonprincipal ultrafilter \(\mathcal U\) on \(\mathbb N\), the ultrapower \({}^\ast\mathbb N=\mathbb N^{\mathbb N}/\mathcal U\) is a proper elementary extension.
We use this framework with the sequence \(p_n\) of prime numbers.

\section{Definitions}
\begin{definition}[Prime sequence]
\label{def:prime-seq}
Let \(p_n\) be the \(n\)-th prime number.
\end{definition}

\begin{definition}[Ultrapower class]
\label{def:ultrapower-class}
Fix a nonprincipal ultrafilter \(\mathcal U\) on \(\mathbb N\).
For a sequence \(a=(a_n)\in\mathbb N^{\mathbb N}\), denote its class in \({}^\ast\mathbb N\) by \([a]_{\mathcal U}\).
\end{definition}

\begin{definition}[\(\Omega_p\)]
\label{def:omega-p}
Define
\[
\Omega_p := [(p_n)]_{\mathcal U}\in {}^\ast\mathbb N.
\]
\end{definition}

\section{Minimal Lemmas}
\begin{lemma}
\label{lem:cofinite-prime-tail}
For every standard \(m\in\mathbb N\), the set \(\{n\in\mathbb N: p_n>m\}\) is cofinite.
\end{lemma}
\begin{proof}
Since \(p_n\to\infty\), there exists \(N\) such that \(n\ge N\Rightarrow p_n>m\).
Hence the set is cofinite.
\end{proof}

\begin{lemma}
\label{lem:ultrafilter-cofinite}
If \(A\subseteq\mathbb N\) is cofinite, then \(A\in\mathcal U\).
\end{lemma}
\begin{proof}
A nonprincipal ultrafilter contains all cofinite sets.
\end{proof}

\section{Main Theorem}
\begin{theorem}[Existence and infinitude of \(\Omega_p\)]
\(\Omega_p\) exists in \({}^\ast\mathbb N\), and for every \(m\in\mathbb N\),
\[
\Omega_p > m
\]
in the ultrapower order.
\end{theorem}
\begin{proof}
Existence is immediate from Definition~\ref{def:omega-p}.
Fix \(m\in\mathbb N\). By Lemma~\ref{lem:cofinite-prime-tail}, \(A_m:=\{n:p_n>m\}\) is cofinite.
By Lemma~\ref{lem:ultrafilter-cofinite}, \(A_m\in\mathcal U\).
By definition of ultrapower order, \([(p_n)]_{\mathcal U}>[(m)]_{\mathcal U}\), i.e. \(\Omega_p>m\).
\end{proof}

\begin{proposition}[Sign-symmetric extension]
Defining \(\Omega_p^-:=-\Omega_p\) in \({}^\ast\mathbb Z\), one has \(\Omega_p^-< -m\) for all \(m\in\mathbb N\).
\end{proposition}
\begin{proof}
Immediate from the theorem and order compatibility with additive inverse.
\end{proof}

\section{What is Proved vs. Conjectural}
\subsection*{Proved in this paper}
\begin{itemize}
\item Existence of \(\Omega_p\) as an ultrapower class.
\item \(\Omega_p\) is larger than every standard natural number.
\item Basic signed extension in \({}^\ast\mathbb Z\).
\end{itemize}

\subsection*{Not proved here (conjectural or deferred)}
\begin{itemize}
\item Any implication for RH or zeros of \(\zeta\).
\item Any identification with wheel-bottom symbols (\(0/0\), \(\infty/\infty\)).
\item Any ``universal residue'' principle across analytic and algebraic settings.
\end{itemize}

\begin{conjecture}[Deferred programmatic claim]
There exists an analytically meaningful extension \(\zeta_{\Omega_p}\) compatible with classical \(\zeta\) on standard inputs and suitable for nonstandard asymptotic analysis.
\end{conjecture}

\section{Limitations}
The construction depends on a nonprincipal ultrafilter (non-constructive existence).
No uniqueness of \(\Omega_p\) across different ultrafilters is claimed.
No classical-analytic theorem is derived from this construction alone.

\section{Conclusion}
We established a conservative foundational step: \(\Omega_p\) exists as an infinite hypernatural tied to the prime sequence.
Further articles must add independent definitions and proofs before any claim on RH or other Millennium problems.

\bibliographystyle{plain}
\begin{thebibliography}{9}
\bibitem{keisler} H. J. Keisler, \emph{Elementary Calculus: An Infinitesimal Approach}.
\bibitem{goldblatt} R. Goldblatt, \emph{Lectures on the Hyperreals}.
\bibitem{marker} D. Marker, \emph{Model Theory: An Introduction}.
\end{thebibliography}

\end{document}
